\documentclass[UTF8,a4paper,12pt]{ctexart}
\usepackage{geometry}
\geometry{margin=2.1cm}
\usepackage{xcolor}
\definecolor{ShanghaiBlue}{HTML}{1F4E79}
\usepackage{longtable}
\usepackage{array}
\usepackage{booktabs}
\usepackage{enumitem}
\usepackage{listings}
\usepackage{fancyhdr}
\usepackage{hyperref}
\hypersetup{colorlinks=true,linkcolor=ShanghaiBlue,urlcolor=ShanghaiBlue}
\setlist[itemize]{leftmargin=2em,itemsep=0.2em,topsep=0.3em}
\setlist[enumerate]{leftmargin=2em,itemsep=0.25em,topsep=0.3em}
\lstset{basicstyle=\ttfamily\small,breaklines=true,columns=fullflexible,frame=single,rulecolor=\color{black!25},backgroundcolor=\color{black!3},xleftmargin=0.5em,xrightmargin=0.5em}
\pagestyle{fancy}
\setlength{\headheight}{15pt}
\fancyhf{}
\lhead{上科大一图通}
\rhead{Techathon 48H}
\cfoot{\thepage}
\begin{document}
\begin{center}\Huge\bfseries\color{ShanghaiBlue}上科大一图通技术文档与第三方资源说明\end{center}
\vspace{0.8em}
\textbf{参赛赛事：}Techathon 48H 黑客马拉松 · 开造 JUST MAKE IT\par
\textbf{学校/单位：}上海科技大学\par
\textbf{提交版本：}2026 年 10 月 18 日\par
\textbf{指导教师：}无\par
\section{1. 技术摘要}
上科大一图通是一个前台地图应用、独立管理后台和 Node.js 服务组成的校园信息系统。前台与后台均使用 Vue 3 + TypeScript + Vite 构建，地图使用 Leaflet；服务端使用 Express 5 与 Node.js 内置 SQLite 接口，静态地图数据和动态业务数据分别存放在两个 SQLite 数据库中。系统支持纯静态地图回退，也可以在配置数据源后运行活动同步、校园账号、社区、排队反馈和 AI 总结服务。\par
运行基线：Node.js 22.13+、npm、现代浏览器；移动端定位需要 HTTPS 或 localhost。主要命令：\par
\begin{lstlisting}
npm install
npm run dev       # 前台/后台开发服务器，默认 5180，/api 代理到 8790
npm run server    # 启动 Express 服务
npm run build     # vue-tsc + Vite 双入口生产构建
npm run start     # 构建后启动服务
npm test          # Vitest 测试
\end{lstlisting}
密钥、SMTP/IMAP 地址和管理员口令只由服务端读取，\texttt{\detokenize{.env}} 不进入前端构建。\par
\section{2. 系统架构与数据流}
\begin{lstlisting}
+------------------------------+
| Vue 前台 / MapView            |
| 搜索、地图、导航、日程、社区  |
+---------------+--------------+
                | REST / JSON / 图片
+---------------v--------------+
| Express 服务                  |
| places · events · canteens    |
| account · schedule · community|
| crowd · checkin · admin       |
+---------------+--------------+
        +-------v--------+  +---v-------------+
        | static.db      |  | dynamic.db      |
        | 建筑/点位/地下 |  | 活动/账号/社区  |
        | 路网/地图配置  |  | 菜单/排队/个人数据|
        +----------------+  +-----------------+
\end{lstlisting}
\texttt{\detokenize{index.html}} 与 \texttt{\detokenize{admin/index.html}} 是两个独立构建入口。前台启动时先加载账号、个人数据、地图和发现目录；服务端接口不可用时，地图和基础导航使用打包在 \texttt{\detokenize{public/data/}} 的静态数据。活动和餐饮目录各自加载，某一来源失败不会阻塞地图。\par
\section{3. 数据基准与规模}
下列数字来自仓库内 \texttt{\detokenize{public/data/}} 的 2026-10-18 快照，用于说明数据结构和演示规模；活动、菜单和后台数据库会随同步与维护变化，不能作为永久库存承诺。\par
\begin{longtable}{>{\raggedright\arraybackslash}p{0.317	extwidth}>{\raggedright\arraybackslash}p{0.317	extwidth}>{\raggedright\arraybackslash}p{0.317	extwidth}}
\toprule
数据 & 当前快照 & 结构/用途 \\
\midrule
\endhead
建筑/区域 & 66 & \texttt{\detokenize{public/data/places.json}} 的 \texttt{\detokenize{buildings}} \\
\midrule
独立点位 & 16 & \texttt{\detokenize{public/data/places.json}} 的 \texttt{\detokenize{pois}} \\
\midrule
房间记录 & 551 & 随建筑保存，用于房间与教室号搜索 \\
\midrule
活动 & 1,022 & \texttt{\detokenize{public/data/events.json}}；官网与邮件来源会动态变化 \\
\midrule
食堂 & 5 & \texttt{\detokenize{public/data/canteens.json}} \\
\midrule
营业点 & 10 & 食堂下的 \texttt{\detokenize{outlets}} \\
\midrule
菜品 & 1,133 & 窗口下的 \texttt{\detokenize{dishes}}，含规格/标签/供应时段 \\
\midrule
地下路网节点 & 75 & \texttt{\detokenize{public/data/underground.json}} version 2 \\
\midrule
地下双向边 & 88 & 每条边只保存两个端点 \\
\midrule
地面/地下连接点 & 26 & \texttt{\detokenize{ramp/stairs/elevator/unknown}} \\
\midrule
地下参考图纸引用 & 6 & \texttt{\detokenize{sources}}；图像文件需与部署包一起核验 \\
\midrule
地下街景点 & 0（静态快照） & 运行时由后台数据库和图片目录维护 \\
\midrule
GPS 校正控制点 & 175 & \texttt{\detokenize{campus.json.warp.src/res}} \\
\midrule
地面栅格 & 250×172 & 每格约 5.4 米，用于步行寻路 \\
\bottomrule
\end{longtable}
\vspace{0.2em}
工作库 \texttt{\detokenize{server/static.db}} 和演示库 \texttt{\detokenize{server/demo.db}} 可以因后台人工维护与演示种子而与 \texttt{\detokenize{public/data/}} 不同；提交演示时应固定使用一套数据快照并在展示说明中注明日期。\par
\section{4. 核心实现}
\subsection{4.1 坐标、底图和地图渲染}
\texttt{\detokenize{src/lib/geo.ts}} 实现旧校园手绘瓦片所需的百度墨卡托 CRS、缩放边界、视野边界和地图平移限制。\texttt{\detokenize{src/lib/geoConvert.ts}} 先把 WGS-84 GPS 转为 GCJ-02，再转为 BD-09，最后利用 \texttt{\detokenize{campus.json}} 的控制点执行仿射变换与残差插值。地图数据以 BD-09 坐标对齐手绘底图，定位结果不会直接把真实地理坐标当作像素坐标。\par
\texttt{\detokenize{MapView.vue}} 使用 Leaflet 图层分离底图、建筑轮廓、点位、活动钉、我的地图旗标、路线、地下层和街景过渡。标签有分类显隐、楼群聚合、选中优先级、可用位置槽、遮挡检测和悬停临时显示；桌面侧栏宽度、手机安全区、底部面板档位和地图可见范围由 \texttt{\detokenize{src/lib/ui.ts}} 统一计算。\par
\subsection{4.2 地面路网与导航}
\texttt{\detokenize{src/lib/route.ts}} 将校园栅格中的可走单元构成八邻接图，按米制步长运行 Dijkstra。搜索缓冲区使用 TypedArray 复用，堆使用可扩容的二叉最小堆；建筑目标支持轮廓、入口或任意点，端点先吸附到可走格。多终点搜索一次处理起点到多个入口的距离，反向距离下界用于剪枝，路线生成阶段用可走性检查将冗余折线拉直。\par
\texttt{\detokenize{src/lib/integratedRoute.ts}} 对地面直达和一次地下通行候选做统一比较。连接点默认在地面端点 45 米内临时投影接入，地下图寻路结果携带有向边序列；集成层缓存连接点对的地下结果，并在当前最优候选已不可能被超越时停止精确计算。实现明确不支持一次路线多次进出地下，也不把几何相交但未显式连接的地下道路自动接通。\par
\subsection{4.3 地下路网、街景和后台编辑}
地下数据采用 \texttt{\detokenize{version: 2}}，核心对象为 \texttt{\detokenize{nodes}}、\texttt{\detokenize{edges}}、\texttt{\detokenize{connectors}}、\texttt{\detokenize{points}} 和 \texttt{\detokenize{sources}}。边是双向的节点对；连接点保存楼层、类型、位置、可选绑定节点、来源和说明；街景点保存两个边端点及插值权重，图片文件按 GUID 存放在 \texttt{\detokenize{data/streetview/}}，不将二进制塞进 JSON 或数据库字段。\par
\texttt{\detokenize{UndergroundWorkspace.vue}} 提供选择、画路、出入口、街景点、插值、删除、吸附、绑定、连通性检查、撤销/重做和显式保存。\texttt{\detokenize{streetViewOrder.ts}} 根据地下路线的有向边序列把边上街景点按行进方向排序；\texttt{\detokenize{streetViewTransition.ts}} 提供街景切换动画，并在减少动态效果设置下缩短转场。\par
\subsection{4.4 个人安排、时间与课表解析}
\texttt{\detokenize{src/lib/today.ts}} 用统一的 \texttt{\detokenize{PlanItem}} 表达课程、活动、吃饭和待办；\texttt{\detokenize{planner.ts}} 计算站点间步行、饭点空档和建议出发时间；\texttt{\detokenize{campusTime.ts}} 统一处理上海时区、活动状态、食堂餐次和演示时间。\texttt{\detokenize{userData.ts}} 将 \texttt{\detokenize{today}}、\texttt{\detokenize{map}}、\texttt{\detokenize{prefs}} 分为可合并的数据项，登录后通过乐观版本号同步，离线时先写本机。\par
\texttt{\detokenize{server/schedule.mjs}} 接收课表图片，服务端调用 DeepSeek Flash 并解析 JSON 结果；前端先展示预览，用户确认后写入个人安排。密钥只在服务端环境变量中读取，接口不把原图长期写入业务数据库。\par
\subsection{4.5 活动数据管道}
\texttt{\detokenize{server/eventSync.mjs}} 调用学校官网活动接口，\texttt{\detokenize{server/mailEvents.mjs}} 解析已授权的通知邮件，\texttt{\detokenize{server/eventPipeline.mjs}} 负责来源调度、游标、去重、状态和待审核队列。活动以稳定 ID、标题、日期、时间和地点共同判断重复；地点解析器支持别名和教室号规则。官网或邮箱失败时保留旧目录；人工修改通过 \texttt{\detokenize{manualOverrides}} 叠加在同步结果上。\par
\texttt{\detokenize{server/organizer.mjs}} 提供主办方活动发布与修改；\texttt{\detokenize{server/checkin.mjs}} 使用活动 ID、时间窗口和服务端密钥生成签到码；管理后台支持审核、置顶、隐藏、取消、讨论管理和数据源运行记录。\par
\subsection{4.6 餐饮与排队}
\texttt{\detokenize{src/lib/menu.ts}} 和 \texttt{\detokenize{canteenMenu.ts}} 负责菜品搜索、价格规格、标签、供应日期、窗口分组和食堂摘要；营业状态由已收录餐次计算。\texttt{\detokenize{server/crowd.mjs}} 对最近 30 分钟的窗口反馈按时间衰减聚合，并执行匿名会话和单窗口 10 分钟限报；没有足够实时反馈时，\texttt{\detokenize{src/lib/crowd.ts}} 回退到基于时段的估算。\par
\subsection{4.7 账号、社区和 AI 总结}
\texttt{\detokenize{server/account.mjs}} 提供邮箱验证码登录、匿名会话绑定、昵称、角色和个人数据同步。会话使用 HttpOnly、SameSite Cookie；登录时可将匿名会话的帖子、互动和个人数据合并到账号。\par
\texttt{\detokenize{server/community.mjs}} 负责话题、帖子、评论、回复、点赞、收藏、图片、菜品评分和 AI 总结。写请求检查同源和 \texttt{\detokenize{X-Community-Request}} 请求头；图片限制格式与大小并清除 EXIF。AI 总结按资料版本和讨论内容指纹缓存，限制并发与调用频率，资料变化后失效。\par
\subsection{4.8 存储、一致性和容错}
\texttt{\detokenize{server/databases.mjs}} 将地图数据放在 \texttt{\detokenize{static.db}}，将活动、菜单、账号、社区、排队、签到和个人数据放在 \texttt{\detokenize{dynamic.db}}，使用 WAL 与迁移逻辑。地图写入先做输入、坐标、几何、GUID 和地下图不变量校验，再在事务中写入；\texttt{\detokenize{surfaceNetworks.mjs}} 在修改栅格前执行在线备份，并对数据库占用返回可读的 503。\par
\section{5. 关键接口摘要}
\begin{longtable}{>{\raggedright\arraybackslash}p{0.475	extwidth}>{\raggedright\arraybackslash}p{0.475	extwidth}}
\toprule
接口 & 用途 \\
\midrule
\endhead
\texttt{\detokenize{GET /api/places}} & 前后台共享的建筑、点位和分类 \\
\midrule
\texttt{\detokenize{GET /api/underground}} / \texttt{\detokenize{PUT /api/underground}} & 读取或保存地下 version 2 图 \\
\midrule
\texttt{\detokenize{GET /api/events}} / \texttt{\detokenize{GET /api/canteens}} & 活动与食堂目录 \\
\midrule
\texttt{\detokenize{POST /api/schedule/parse}} & 课表截图解析 \\
\midrule
\texttt{\detokenize{POST /api/account/code}} / \texttt{\detokenize{verify}} & 邮箱验证码登录 \\
\midrule
\texttt{\detokenize{GET /api/me/data}} / \texttt{\detokenize{PUT /api/me/data/:key}} & 登录后的个人数据同步 \\
\midrule
\texttt{\detokenize{GET/POST /api/organizer/events}} & 主办方活动发布 \\
\midrule
\texttt{\detokenize{GET/POST /api/events/pipeline}} & 活动数据源状态与运行 \\
\midrule
\texttt{\detokenize{GET/POST /api/events/review}} & 活动待审核队列 \\
\midrule
\texttt{\detokenize{GET/POST /api/crowd}} & 排队聚合与匿名反馈 \\
\midrule
\texttt{\detokenize{/api/community/*}} & 话题、帖子、评论、互动、图片和总结 \\
\midrule
\texttt{\detokenize{/api/admin/*}} & 管理员活动、菜单、社区和系统操作 \\
\bottomrule
\end{longtable}
\vspace{0.2em}
\section{6. 安全、隐私与授权边界}
\begin{itemize}
\item \texttt{\detokenize{.env}}、管理员口令、签到密钥、SMTP/IMAP 凭据和 DeepSeek 密钥只在服务端使用，不进入前端 bundle。
\item 社区浏览可匿名；发布、评论、上传和签到需要账号与昵称。匿名会话只用于站内互动与排队反馈，不把会话 token 暴露给帖子作者。
\item 课表截图和个人安排属于个人数据；演示与截图应使用演示账号和脱敏内容，不应把真实邮箱、验证码、教师姓名或用户照片放入提交附件。
\item 活动邮件只抽取标题、时间、地点和主办等必要字段；原始邮件不随代码仓库分发。
\item 菜单照片、活动内容和地下图纸按数据包记录的来源、拍摄者/提供者及比赛授权口径列明；外部素材不在原创声明中表述为团队原创。
\end{itemize}
\section{7. 第三方软件、模型、API 与数据资源}
\subsection{软件依赖}
\begin{longtable}{>{\raggedright\arraybackslash}p{0.237	extwidth}>{\raggedright\arraybackslash}p{0.237	extwidth}>{\raggedright\arraybackslash}p{0.237	extwidth}>{\raggedright\arraybackslash}p{0.237	extwidth}}
\toprule
依赖 & 版本 & 用途 & 许可证 \\
\midrule
\endhead
Vue & 3.5.43 & 前台与后台组件、响应式状态 & MIT \\
\midrule
Leaflet & 1.9.4 & 地图、图层、交互与 CRS 承载 & BSD-2-Clause \\
\midrule
Express & 5.2.1 & HTTP API 与静态服务 & MIT \\
\midrule
Vite & 8.3.1 & 双入口开发与构建 & MIT \\
\midrule
Vitest & 4.1.12 & 单元与接口测试 & MIT \\
\midrule
TypeScript & 5.9.3 & 类型检查与编译 & Apache-2.0 \\
\midrule
imapflow & 2.2.1 & IMAP 通知邮件读取 & MIT \\
\midrule
mailparser & 3.9.33 & 邮件解析 & MIT \\
\midrule
nodemailer & 8.0.11 & 邮箱验证码发送 & MIT \\
\midrule
qrcode-generator & 2.0.4 & 分享与活动二维码生成 & MIT \\
\bottomrule
\end{longtable}
\vspace{0.2em}
版本以 \texttt{\detokenize{package-lock.json}} 为准；间接依赖不在此表逐项展开。\par
\subsection{数据、服务和图片}
\begin{longtable}{>{\raggedright\arraybackslash}p{0.317	extwidth}>{\raggedright\arraybackslash}p{0.317	extwidth}>{\raggedright\arraybackslash}p{0.317	extwidth}}
\toprule
资源 & 用途 & 来源与提交说明 \\
\midrule
\endhead
校园手绘底图、楼宇轮廓、房间表、服务点 & 地图、搜索、入口与定位对齐 & 上海科技大学校园地图 \texttt{\detokenize{map.shanghaitech.edu.cn}}；按赛前数据包的来源说明署名并按比赛提交授权口径说明 \\
\midrule
部分道路与周边地理数据 & 补充步行栅格与地点信息 & OpenStreetMap contributors；按 ODbL 保留署名与相应说明 \\
\midrule
官网活动日历 & 活动目录同步 & 学校官网公开活动接口；仅抽取活动展示所需字段 \\
\midrule
学校通知邮件 & 活动补充来源 & 使用授权通知邮件，抽取标题、时间、地点和主办，不分发原始邮件 \\
\midrule
菜单、营业时间、菜单原图 & 食堂与菜品功能 & 赛前数据包记录的窗口价目牌实拍和公共服务处通知，提交材料注明采集日期与来源 \\
\midrule
菜品实拍 & 菜品详情 & 数据包记录来源为 QQ 空间相册“科大美食”；按授权口径标注，不作为团队原创照片 \\
\midrule
11–13 号楼外观与施工边界 & 补画新楼与在建区域 & 现场照片、施工公示平面图及团队整理；按比赛提交口径说明 \\
\midrule
DeepSeek Chat API / DeepSeek Flash & 课表截图结构化解析、社区资料总结 & 外部模型服务，密钥由服务端管理；模型输出是辅助结果，需由用户确认，不宣称为原创数据 \\
\midrule
百度坐标/墨卡托公开常数 & 坐标转换与底图 CRS & 公开技术常数与公式，代码内按项目数据坐标约定使用 \\
\midrule
\texttt{\detokenize{skdmap}} 原型 & 地图 CRS、步行栅格、拖动限制、界面与图标的实现参考 & 团队自有原型；提交材料应说明复用、重构和新增部分 \\
\bottomrule
\end{longtable}
\vspace{0.2em}
\section{8. 测试与验证}
当前仓库包含 56 个测试文件（\texttt{\detokenize{src}} 41、\texttt{\detokenize{server}} 9、\texttt{\detokenize{scripts}} 6），覆盖路线、坐标、地下图校验与编辑、标签避让、活动解析与同步、菜单、排队、账号、社区、数据库迁移和管理 API。\par
已经执行并通过：\par
\begin{itemize}
\item \texttt{\detokenize{npm run build}}：\texttt{\detokenize{vue-tsc -b}} 与 Vite 双入口构建成功。
\item 选定的核心纯逻辑测试（路线、地面/地下集成、地下图、坐标转换、菜单、排队、规划、时间等）：25 个测试文件、268 个测试通过。
\end{itemize}
完整 \texttt{\detokenize{npx vitest run --exclude '.claude/**'}} 在当前受限运行环境中无法作为全绿证据：服务器测试尝试监听本地端口时收到 \texttt{\detokenize{EPERM}}，引发账号、社区、排队、活动管道和管理员接口测试失败；另有两条 \texttt{\detokenize{mapHighlights}} 测试的固定期望与现行演示数据不一致。该结果反映当前执行环境和测试夹具状态，不能据此否定已通过的纯逻辑测试；在可监听本地回环端口的干净工作区应重新运行服务器测试并修正数据夹具。\par
\section{9. 已知限制}
\begin{enumerate}
\item 当前地下模型固定为一个 \texttt{\detokenize{floor = -1}} 层级，路线只支持一次进入和一次离开地下。
\item 地下连接点、路网和街景照片需要现场核验；静态快照中 \texttt{\detokenize{points}} 为空，街景资源与地下图纸引用需在部署包中检查。
\item \texttt{\detokenize{demoIndoor503.ts}} 是固定演示路线，不能代表任意建筑的通用室内导航。
\item 避雨模式使用独立栅格和当前测试版的强制地下路线策略，仍应把它称为测试/实验能力。
\item 活动、菜单、营业状态和排队估算具有时间性，不代表实时库存、实时营业或现场秩序。
\item 定位受浏览器权限、HTTPS、设备 GPS 和室内信号影响；地下没有 GPS 时依赖街景或人工操作。
\item AI 解析与总结依赖外部服务额度和密钥，返回结果需要用户确认；服务不可用时系统应保留原始资料并给出失败提示。
\item 外部地图、图片、活动、邮件和模型资源在提交材料中按清单列明来源和授权口径。
\end{enumerate}
\section{10. 运行与交付注意事项}
演示建议使用 \texttt{\detokenize{server/demo.db}} 与 \texttt{\detokenize{DEMO_MODE=true}}，使用假账号、脱敏活动和模拟定位；不要把 \texttt{\detokenize{.env}}、\texttt{\detokenize{server/.admin-password}}、\texttt{\detokenize{server/.checkin-secret}}、验证码、真实用户数据或未获授权的原始邮件放入提交包。部署时同时备份 \texttt{\detokenize{server/static.db}}、\texttt{\detokenize{server/dynamic.db}} 和 \texttt{\detokenize{data/streetview/}}，运行中的 SQLite 使用一致性备份方式，避免只复制单个 WAL 文件。\par
\end{document}
